% TEMPLATE-MILC-v3.tex - plantilla maestra UNICA de las guias MILC v3.
%
% La usan las tres familias de guias, cada una con su driver:
%   · Grados 6-11  -> scripts/build-guias-g11.py
%   · Bebras       -> scripts/build-guias-bebras.py
%   · Semillero    -> scripts/build-guias-semillero.py
%
% Antes habia tres copias casi identicas (~400 lineas iguales, 6 distintas):
% cada mejora habia que aplicarla tres veces, y en la practica no se hacia
% (el motor de recursos vivia solo en dos de ellas). Lo que varia entre
% familias esta parametrizado en 6 placeholders que cada driver llena
% (se nombran aqui SIN su sintaxis real: el builder reemplaza por texto plano,
% tambien dentro de los comentarios, y un valor multilinea rompe el preambulo):
%
%   HEADER_IZQ        encabezado izquierdo de pagina
%   HEADER_DER        encabezado derecho de pagina
%   PORTADA_ETIQUETA  rotulo sobre el numero grande (GRADO/MOMENTO/MODULO)
%   PORTADA_SERIE     linea de serie de la portada
%   PORTADA_ID        identificador de la guia en la portada
%   PORTADA_CIRCULO   el circulito de grado (°), o vacio si no aplica
%
% El resto de claves las documenta content/guias/_SCHEMA.md. El script valida
% que no queden placeholders sin reemplazar antes de escribir el archivo final.

\documentclass[11pt,letterpaper]{article}
\usepackage[margin=1.75cm]{geometry}
\usepackage[table]{xcolor}
\usepackage{fontspec}
\usepackage[spanish]{babel}
\usepackage{tikz}
% Librerías TikZ para los diagramas de las guías (bloque `recursos:`):
%  · babel  → babel en español vuelve activos caracteres como «>», y TikZ los usa
%             en sus opciones (>=stealth, ->). Sin esta librería, cualquier
%             diagrama con flechas revienta con «\language@active@arg> has an
%             extra }». Debe ir ANTES de que se dibuje nada.
%  · positioning        → sintaxis «below left=2pt and 3pt».
%  · decorations.pathreplacing → llaves ({brace}) para señalar rangos.
\usetikzlibrary{babel,positioning,decorations.pathreplacing}
\usepackage{graphicx}
% Circuitos: el trabajo de electrónica se hace hoy con micro:bit y sus sensores
% integrados (MakeCode, sin cableado), así que no se necesitan esquemáticos.
% Si algún día entran montajes con protoboard, basta descomentar esta línea:
% los diagramas .tex/.tikz declarados en `recursos:` ya se incrustan solos.
% \usepackage{circuitikz}
\usepackage[most]{tcolorbox}
\usepackage{tabularx}
\usepackage{array}
\usepackage{fancyhdr}
\usepackage{titlesec}
\usepackage{ragged2e}
\usepackage{enumitem}
\usepackage{microtype}

% Helvetica Neue no trae la flecha U+2192, asi que cada «→» escrito literalmente
% en el YAML se compilaba a NADA, en silencio (462 flechas en 61 guias: rutas del
% tipo «paleta → tipografia → lockup» salian rotas). Mapeamos el caracter a su
% version matematica, que si tiene glifo. Asi el autor puede seguir escribiendo
% «→» comodamente en el YAML.
\usepackage{newunicodechar}
\newunicodechar{→}{\ensuremath{\rightarrow}}

\setmainfont{Helvetica Neue}
\setsansfont{Helvetica Neue}
\setlength{\parindent}{0pt}
\setlength{\parskip}{6pt}
\hyphenpenalty=200
\exhyphenpenalty=200
\pretolerance=400
\tolerance=2000
\setlength{\emergencystretch}{1em}
\renewcommand{\arraystretch}{1.25}
\setlist{leftmargin=5mm,itemsep=1mm,topsep=1mm}
\newcolumntype{Y}{>{\RaggedRight\arraybackslash}X}

\definecolor{milcMagenta}{HTML}{E6007E}
\definecolor{milcVino}{HTML}{5A0038}
\definecolor{milcTurquesa}{HTML}{0093A5}
\definecolor{milcVerde}{HTML}{58A923}
\definecolor{milcMostaza}{HTML}{E5B400}
\definecolor{milcOcre}{HTML}{B86600}
\definecolor{milcNegro}{HTML}{241816}
\definecolor{milcGris}{HTML}{F7F7F4}
\definecolor{milcLinea}{HTML}{E8E2DD}
\definecolor{milcGradePrimary}{HTML}{84CC16}
\definecolor{milcGradeDark}{HTML}{4D7C0F}
\definecolor{milcGradeSoft}{HTML}{ECFCCB}
\definecolor{milcGradeAccent}{HTML}{CFE7FF}
\definecolor{milcGradeText}{HTML}{FFFFFF}
\color{milcNegro}

\pagestyle{fancy}
\fancyhf{}
\lhead{\small Tecnología e Informática · Grado 7}
\rhead{\small Guía 21 · MILC v3}
\cfoot{\small\thepage}
\renewcommand{\headrulewidth}{0pt}

\newtcolorbox{titlebox}[1]{
  enhanced,colback=#1,colframe=#1,arc=7mm,boxrule=0pt,
  left=7mm,right=7mm,top=3mm,bottom=3mm,
  before skip=8pt,after skip=8pt,
  fontupper=\sffamily\bfseries\Large\color{white}
}

\newtcolorbox{softbox}[3]{
  enhanced,colback=#2,colframe=#1,borderline west={4pt}{0pt}{#1},
  arc=2mm,boxrule=.4pt,left=4mm,right=4mm,top=3mm,bottom=3mm,
  before skip=6pt,after skip=8pt,title={#3},coltitle=white,
  colbacktitle=#1,fonttitle=\sffamily\bfseries,fontupper=\RaggedRight,
  attach boxed title to top left={xshift=3mm,yshift=-2mm},
  boxed title style={arc=2mm, boxrule=0pt}
}

\newtcolorbox{drawbox}[2]{
  enhanced,colback=white,colframe=#1,arc=4mm,boxrule=1pt,
  left=5mm,right=5mm,top=4mm,bottom=4mm,before skip=8pt,after skip=8pt,
  title={#2},coltitle=white,colbacktitle=#1,fonttitle=\sffamily\bfseries,
  fontupper=\RaggedRight,
  attach boxed title to top left={xshift=4mm,yshift=-2mm},
  boxed title style={arc=3mm, boxrule=0pt}
}

\newtcolorbox{infoband}[1]{
  enhanced,colback=#1!8,colframe=#1!50,arc=2mm,boxrule=.5pt,
  left=4mm,right=4mm,top=2mm,bottom=2mm,before skip=4pt,after skip=4pt,
  fontupper=\small\RaggedRight
}

% Puente narrativo entre secciones (contrato editorial Conéctate).
% Da continuidad: "Acabas de X. Ahora vas a Y porque Z."
% Visualmente: banda izquierda en color del grado, texto en itálica pequeña.
\newtcolorbox{puentebox}{
  enhanced,colback=milcGradeSoft,colframe=milcGradeDark,
  borderline west={3pt}{0pt}{milcGradeDark},
  arc=0mm,boxrule=0pt,left=5mm,right=4mm,top=2.5mm,bottom=2.5mm,
  before skip=4pt,after skip=8pt,
  fontupper=\itshape\small\color{milcNegro}\RaggedRight
}
% La flecha va en modo matematico a proposito: Helvetica Neue NO tiene el glifo
% U+2192, asi que \textrightarrow se compilaba a NADA («Missing character: There
% is no → in font Helvetica Neue») y el puente salia sin flecha. En modo
% matematico la flecha viene de la fuente matematica y si se dibuja.
\newcommand{\puente}[1]{%
  \begin{puentebox}%
    \textcolor{milcGradeDark}{$\rightarrow$}\hspace{2mm}#1%
  \end{puentebox}%
}

\newcommand{\linead}[1]{\noindent\color{milcLinea}\rule{#1}{0.5pt}\color{milcNegro}}
\newcommand{\respuesta}[1]{\par\vspace{2mm}\foreach \n in {1,...,#1}{\noindent\color{milcLinea}\rule{\linewidth}{0.5pt}\par\vspace{5mm}}\color{milcNegro}}
\newcommand{\checkbox}{\textcolor{milcGradeDark}{\fbox{\phantom{x}}}\hspace{5pt}}

\newcommand{\phasecard}[4]{
\begin{tcolorbox}[enhanced,colback=#3,colframe=#2,arc=3mm,boxrule=.5pt,height=3.05cm,valign=center,halign=center,left=1.5mm,right=1.5mm,top=2mm,bottom=2mm]
{\sffamily\bfseries\fontsize{22}{24}\selectfont\textcolor{#2}{#1}}\par
{\sffamily\bfseries\small #4}
\end{tcolorbox}
}

% ── Recursos visuales de la guía ──────────────────────────────────────
% Declarados en el bloque `recursos:` del YAML y emitidos por el builder.
% \guiaFigura   → imagen rasterizada o PDF (foto, carta celeste, montaje).
% \guiaDiagrama → fuente TikZ/circuitikz (.tex) incrustada como vector:
%                 es la vía para circuitos y esquemáticos editables.
% #1 = ruta relativa al .tex · #2 = pie (puede ir vacío)
\newcommand{\guiaFigura}[2]{%
\begin{center}
\includegraphics[width=0.88\linewidth,height=0.38\textheight,keepaspectratio]{#1}\\[1.5mm]
{\footnotesize\itshape\textcolor{milcNegro!75}{#2}}
\end{center}
}

\newcommand{\guiaDiagrama}[2]{%
\begin{center}
\input{#1}\\[1.5mm]
{\footnotesize\itshape\textcolor{milcNegro!75}{#2}}
\end{center}
}

\begin{document}
\thispagestyle{empty}
\begin{tikzpicture}[remember picture,overlay]
  \fill[white] (current page.south west) rectangle (current page.north east);
  \path[fill=milcGradePrimary,rounded corners=16mm]
    ([xshift=.55cm,yshift=-.55cm]current page.north west) rectangle
    ([xshift=-.55cm,yshift=3.65cm]current page.south east);

  \node[anchor=north west,text=milcGradeText] at ([xshift=2.0cm,yshift=-1.85cm]current page.north west)
    {\sffamily\bfseries\fontsize{14}{16}\selectfont GRADO};
  \node[anchor=north west,text=milcGradeText,opacity=.82] at ([xshift=2.0cm,yshift=-2.75cm]current page.north west)
    {\sffamily\bfseries\fontsize{9}{11}\selectfont SERIE GUÍAS · TECNOLOGÍA E INFORMÁTICA};

  \begin{scope}[shift={([xshift=-3.05cm,yshift=-2.55cm]current page.north east)}]
    \fill[white,opacity=.80] (-.62,-.52) rectangle (.62,.52);
    \draw[milcGradeText,line width=1pt] (-.62,-.52) rectangle (.62,.52);
    \fill[milcGradeDark] (-.42,-.38) rectangle (-.22,.06);
    \fill[milcGradeAccent] (-.10,-.38) rectangle (.10,.24);
    \fill[milcGradePrimary!60!black] (.22,-.38) rectangle (.42,.38);
  \end{scope}

  \node[anchor=west,text=milcGradeText] at ([xshift=1.9cm,yshift=-12.85cm]current page.north west)
    {\sffamily\bfseries\fontsize{124}{124}\selectfont 7};
  % El circulito de grado (°) solo aplica a las guias de grado («11°»); en
  % Bebras y Semillero el numero grande es un momento/modulo, no un grado.
  \draw[milcGradeText,line width=1.7pt]
    ([xshift=4.52cm,yshift=-11.68cm]current page.north west) circle (.13cm);

  \node[anchor=west,text=milcGradeText,text width=8.7cm,align=left] at ([xshift=10.35cm,yshift=-11.6cm]current page.north west)
    {
     {\sffamily\bfseries\fontsize{19}{22}\selectfont Apertura\\periodo 3 ---\\la IA\\como nuevo\\consejero}\\[4mm]
     {\sffamily\bfseries\fontsize{23}{26}\selectfont Guía 21-7-TIC}\\[2mm]
     {\sffamily\fontsize{13}{16}\selectfont Período 3 · Inteligencia Artificial}\\[5mm]
     {\sffamily\bfseries\fontsize{11}{13}\selectfont Institución Educativa Sor María Juliana}\\[1mm]
     {\sffamily\fontsize{10}{12}\selectfont Autor: Dr. Álvaro Cárdenas Orozco}
    };

  \path[fill=milcGradeSoft,rounded corners=7mm]
    ([xshift=1.35cm,yshift=.85cm]current page.south west) rectangle
    ([xshift=-1.35cm,yshift=4.25cm]current page.south east);
  \draw[milcGradeDark,line width=.65pt,rounded corners=7mm]
    ([xshift=1.35cm,yshift=.85cm]current page.south west) rectangle
    ([xshift=-1.35cm,yshift=4.25cm]current page.south east);
  \node[anchor=center,text=milcNegro,text width=17.0cm,align=center] at ([yshift=2.55cm]current page.south)
    {\sffamily\fontsize{10}{12}\selectfont Metodología MILC · Apertura ancestral · Pensamiento computacional · Triángulo Dussel-Estoicismo-Floridi\\
    \textcolor{milcGradeDark}{\bfseries Producto:} Tu \textbf{cuaderno listo para el periodo 3 de grado 7}: \textbf{índice del P3} con las 10 sesiones + \textbf{5 acuerdos del uso responsable de IA} firmados (verificar, no compartir datos personales, citar la fuente, pensar antes de copiar, mantener tu voz propia) + \textbf{5 preguntas-radar} sobre IA + \textbf{mini-investigación}: \emph{``¿Quién era el consejero del barrio en la generación de mis abuelos? ¿Qué hacía? ¿Qué consejos daba?''} (pregunta a familiares mayores).};
\end{tikzpicture}
\mbox{}
\newpage

\section*{Apertura: Apertura periodo 3 --- la IA como el nuevo consejero del barrio}

\begin{softbox}{milcGradeDark}{milcGradeSoft}{Saber ancestral}
\textbf{En el barrio de Cartago, cada cuadra tenía su consejero.} No era un cargo formal --- era \emph{la persona mayor a quien todos acudían cuando había una decisión difícil}. Doña Mercedes era consejera para temas de niños y educación. Don Lucho el relojero era consejero para temas técnicos. Doña Esperanza la partera era consejera para temas de salud familiar. \textbf{¿Por qué la gente confiaba en ellos?} Porque tenían \emph{experiencia} (habían visto muchas vidas), \emph{paciencia} (escuchaban antes de aconsejar), \emph{honestidad} (no fingían saber lo que no sabían), \emph{discreción} (lo que se les contaba no se difundía). La gente preguntaba; el consejero respondía; la gente \textbf{verificaba con su propio criterio} y decidía. \emph{El consejero no decidía por ti; te ayudaba a decidir mejor}. Hoy hay un nuevo tipo de consejero al que millones de personas acuden cada día: la \textbf{Inteligencia Artificial} (ChatGPT, Claude, Gemini). Te puede aconsejar sobre cualquier tema en segundos. Pero con la IA pasa lo mismo que con doña Mercedes: \emph{si preguntas bien, verificas con criterio y mantienes tu voz, te sirve. Si copias sin pensar, te perjudica}.
\end{softbox}

\begin{softbox}{milcTurquesa}{milcGris}{Saber contemporáneo}
Este periodo se llama \textbf{Inteligencia Artificial}. \emph{Estás aprendiendo IA en el momento histórico} en que esta tecnología está cambiando el mundo. Tu generación es la primera que crece con IA conversacional accesible (ChatGPT salió a fines de 2022). En 10 sesiones vas a recorrer: \textbf{(S1) Hoy: Apertura --- la IA como consejero.} \textbf{(S2) ¿Qué es la IA?} \textbf{(S3) Tipos de IA.} \textbf{(S4) Cómo aprende una IA (datos y sesgo).} \textbf{(S5) Modelos de lenguaje (ChatGPT, Claude, Gemini).} \textbf{(S6) Prompting básico.} \textbf{(S7) Prompting avanzado.} \textbf{(S8) Ética de la IA.} \textbf{(S9) IA y mis tareas.} \textbf{(S10) Cosecha: proyecto con IA responsable.} \textbf{Aprenderás a usar IA con criterio:} preguntar bien, verificar respuestas, citar cuando uses ideas, mantener tu voz propia. \emph{No es tema de futuro: es tema de hoy. Las decisiones que tomes con IA en grado 7 forman el hábito que llevarás al resto de tu vida}.
\end{softbox}

\begin{softbox}{milcMostaza}{milcGris}{Pregunta-puente}
\textit{Si tu mamá te pregunta \emph{``¿qué es ChatGPT?''}, ¿\textbf{podrías explicárselo} con tus palabras? ¿Y si te preguntara si es buena idea que ella lo use para escribir un correo profesional?}
\end{softbox}

\begin{softbox}{milcVerde}{milcGris}{Saber y saber hacer}
Al terminar la clase: (1) podrás \textbf{identificar} la ruta del periodo 3; (2) sabrás \textbf{explicar} por qué la IA es como un consejero moderno; (3) podrás \textbf{aplicar} los 5 acuerdos del uso responsable; (4) habrás \textbf{creado} tu cuaderno P3 + investigación sobre consejeros tradicionales.
\end{softbox}



\small
\begin{tabularx}{\textwidth}{>{\bfseries}p{3.0cm}Y}
\rowcolor{milcGradePrimary!12}Autor & Dr. Álvaro Cárdenas Orozco\\
DBA & Reconoce los fundamentos, usos y límites éticos de la Inteligencia Artificial (MEN, T\&I 7°)\\
Referentes & Saber ancestral del consejero · Modelos de lenguaje · Ética algorítmica y prompting\\
Producto final & Tu \textbf{cuaderno listo para el periodo 3 de grado 7}: \textbf{índice del P3} con las 10 sesiones + \textbf{5 acuerdos del uso responsable de IA} firmados (verificar, no compartir datos personales, citar la fuente, pensar antes de copiar, mantener tu voz propia) + \textbf{5 preguntas-radar} sobre IA + \textbf{mini-investigación}: \emph{``¿Quién era el consejero del barrio en la generación de mis abuelos? ¿Qué hacía? ¿Qué consejos daba?''} (pregunta a familiares mayores).\\
\end{tabularx}
\normalsize

\puente{Hoy haces 4 cosas: descubres la analogía consejero/IA, conoces la ruta del periodo, firmas los acuerdos, planeas investigación.}

\begin{titlebox}{milcGradeDark}
Ruta MILC: así vamos a aprender
\end{titlebox}
\begin{tabularx}{\textwidth}{>{\centering\arraybackslash}X>{\centering\arraybackslash}X>{\centering\arraybackslash}X>{\centering\arraybackslash}X}
\phasecard{1}{milcVerde}{milcGris}{Escucha\\[-1mm]\small Observo, escucho y pregunto.} &
\phasecard{2}{milcTurquesa}{milcGris}{Sistematización\\[-1mm]\small Pensamiento computacional.} &
\phasecard{3}{milcMagenta}{milcGris}{Praxis\\[-1mm]\small Construyo y aplico.} &
\phasecard{4}{milcVino}{milcGris}{Evaluación\\[-1mm]\small Reflexión liberadora.}
\end{tabularx}


\puente{Empezamos imaginando la conversación entre tu abuela y un consejero del barrio.}

\begin{titlebox}{milcVerde}
Fase 1 · Escucha
\end{titlebox}

\begin{softbox}{milcVerde}{milcGris}{Escena para escuchar}
\textbf{Actividad 1 · IDENTIFICA --- La conversación de tu abuela} (12 min · individual). Cierra los ojos e imagina esta escena: \emph{tu abuela, hace 30 años, tiene una decisión difícil --- ¿vende la finca o no la vende? Va donde el consejero del barrio (don Lucho, doña Mercedes, alguien parecido). Le cuenta el dilema. El consejero escucha, hace preguntas, le da una opinión, le advierte qué cosas considerar}. \textbf{En el cuaderno responde:} \emph{(1) ¿Qué cualidades tenía el consejero que hacía que la abuela confiara en él?} \emph{(2) ¿La abuela seguía exactamente el consejo o lo combinaba con su propio criterio?} \emph{(3) ¿Qué pasaba si el consejero se equivocaba? (¿lo culpaba la abuela?)}. \emph{(4) ¿Por qué la abuela buscaba consejo en lugar de decidir sola?}.
\end{softbox}

\checkbox Imaginé la escena con calma.\\
\checkbox Respondí las 4 preguntas con honestidad.\\
\checkbox Pensé en las cualidades del consejero antiguo.

\begin{infoband}{milcVerde}


\textbf{Lo que va al cuaderno (Actividad 1 · IDENTIFICA).} Título: «Actividad 1 --- El consejero del barrio». \textbf{Formato:} 4 respuestas cortas. \textbf{Extensión:} 5-6 líneas. \textbf{Sabes que terminaste cuando} reconoces que el consejero es modelo + tú decides.
\end{infoband}

\puente{Después aprendes la ruta del periodo + 5 acuerdos del uso responsable de IA.}

\begin{titlebox}{milcTurquesa}
Fase 2 · Sistematización con pensamiento computacional
\end{titlebox}

\begin{softbox}{milcTurquesa}{milcGris}{Cuatro pilares aplicados al tema}
Las cualidades del consejero (experiencia, paciencia, honestidad, discreción) son \textbf{exactamente las que esperarías de una IA ideal}. Veamos en qué cumple y en qué no. Después aprendes la ruta + 5 acuerdos.
\end{softbox}

\small
\begin{tabularx}{\textwidth}{>{\bfseries\RaggedRight\arraybackslash}p{3.6cm}Y}
\rowcolor{milcTurquesa!90}\color{white}Pilar & \color{white}\bfseries Aplicación al tema\\
1. Descomposición & \textbf{¿Cómo se parece la IA al consejero del barrio?} \textbf{(a) Experiencia:} la IA aprendió de millones de textos (libros, artículos, foros). En cierto sentido, \emph{ha visto muchas vidas} (las que están en sus datos). \textbf{(b) Paciencia:} no se cansa de responder. Le preguntas 100 veces lo mismo y responde 100 veces. \textbf{(c) Discreción:} \emph{depende mucho del servicio}. Algunos guardan tus conversaciones; otros no. La discreción de la IA NO es como la de un humano: hay servidores de por medio.\\
2. Reconocimiento de patrones & \textbf{¿En qué NO se parece al consejero del barrio?} \textbf{(a) Honestidad:} la IA \emph{a veces inventa cosas con seguridad} (se llama \emph{``alucinación''}). Don Lucho diría \emph{``no sé''}; la IA puede inventar una respuesta plausible que es falsa. \textbf{(b) Conocimiento de contexto:} el consejero conocía a tu familia, tu barrio, tu historia. La IA no conoce nada de eso. Le falta el contexto humano. \textbf{(c) Responsabilidad:} si el consejero del barrio te aconseja mal, la comunidad lo nota y su reputación cae. \emph{La IA no tiene reputación social que perder}. Te aconseja mal y no le pasa nada.\\
3. Abstracción & \textbf{La ruta del P3 --- 10 sesiones.} \emph{S1: Hoy: Apertura.} \emph{S2: ¿Qué es la IA?} \emph{S3: Tipos de IA.} \emph{S4: Cómo aprende una IA (datos y sesgo).} \emph{S5: Modelos de lenguaje.} \emph{S6: Prompting básico.} \emph{S7: Prompting avanzado.} \emph{S8: Ética.} \emph{S9: IA en mis tareas.} \emph{S10: Cosecha.} El periodo se organiza en 3 bloques: \emph{conocer la IA (S2-S5), usarla bien (S6-S7), usarla con ética (S8-S10)}.\\
4. Algoritmo conceptual & \textbf{Los 5 acuerdos del uso responsable de IA.} (1) \textbf{Verifica siempre.} La IA puede equivocarse o inventar. Lo que te diga, lo verificas en al menos 1 fuente humana o de criterio. (2) \textbf{No compartas datos personales sensibles.} Tu dirección, tu cédula, tu contraseña, info íntima. Una vez que escribes algo en el chat, no sabes dónde queda. (3) \textbf{Cita la fuente cuando uses ideas.} Si la IA te ayudó a desarrollar una idea, dilo en tu trabajo (\emph{``con apoyo de Claude AI''}). Es honestidad académica. (4) \textbf{Piensa antes de copiar.} La IA es asistente, no reemplazo. Si copias sin entender, te dañas tú: no aprendes, no creces, te haces dependiente. (5) \textbf{Mantén tu voz propia.} El estilo, la perspectiva, las decisiones siguen siendo tuyas. La IA puede ayudar pero la voz final es tuya.\\
\end{tabularx}
\normalsize

\begin{drawbox}{milcTurquesa}{Ruta P3 + 5 acuerdos}
\textbf{Ruta P3:} \\[1mm]
Bloque 1 (S2-S5): conocer la IA. \\[1mm]
Bloque 2 (S6-S7): usarla bien con prompting. \\[1mm]
Bloque 3 (S8-S10): usarla con ética + cosecha. \\[1mm]
\textbf{5 acuerdos:} verificar · no datos personales · citar · pensar antes de copiar · voz propia.
\end{drawbox}

\begin{softbox}{milcMostaza}{milcGris}{Errores comunes y su corrección}
\textbf{Mitos sobre la IA que escucharás.} (1) \emph{``La IA siempre dice la verdad''} --- Falso. Puede inventar con seguridad (alucinación). Verifica. (2) \emph{``Usar IA es hacer trampa''} --- Depende. Usarla como asistente con atribución no es trampa; copiar sin entender sí. (3) \emph{``La IA va a reemplazar a los humanos''} --- Falso en muchos sentidos. Algunas tareas sí; otras no. La IA aumenta humanos, no los reemplaza completos. (4) \emph{``La IA solo es para programadores''} --- Falso. Hoy cualquier persona la usa: escritores, profesores, médicos, estudiantes. (5) \emph{``Lo que escriba a la IA queda solo conmigo''} --- Falso. Tus conversaciones quedan en servidores. No compartas info sensible.
\end{softbox}

\begin{infoband}{milcTurquesa}


\textbf{Lo que va al cuaderno (Actividad 2 · EXPLICA).} Título: «Actividad 2 --- IA como consejero + ruta P3 + 5 acuerdos». \textbf{Formato:} bloque A con tabla comparativa de \emph{consejero del barrio} vs \emph{IA} (3 similitudes + 3 diferencias). Bloque B con ruta del P3 (10 sesiones en 3 bloques). Bloque C con 5 acuerdos firmados. \textbf{Extensión:} 15 líneas. \textbf{Sabes que terminaste cuando} podrías explicarle a tu mamá qué es la IA usando la analogía del consejero.
\end{infoband}

\puente{Con todo claro, dejas el cuaderno listo + 5 preguntas + investigación planificada.}

\begin{titlebox}{milcMagenta}
Fase 3 · Praxis con pensamiento computacional
\end{titlebox}

\begin{softbox}{milcMagenta}{milcGris}{Construyendo el producto con los cuatro pilares}
\textbf{Actividad 3 · CREA --- Cuaderno P3 + 5 preguntas + investigación del consejero} (28 min · individual + tarea en casa). Vas a dejar el cuaderno listo + plantear preguntas + investigar.
\end{softbox}

\small
\begin{tabularx}{\textwidth}{>{\bfseries\RaggedRight\arraybackslash}p{3.6cm}Y}
\rowcolor{milcMagenta!90}\color{white}Pilar & \color{white}\bfseries Aplicación al producto\\
Descomposición & \textbf{Paso 1 --- Índice del P3 + encabezado fijo.} En tu cuaderno de Tecnología, abre una nueva página. Título: \textbf{``Periodo 3 --- Inteligencia Artificial''}. Debajo, índice de 10 sesiones del P3 (de la ruta de la sistematización).\\
Patrones & \textbf{Paso 2 --- 5 acuerdos firmados.} En la siguiente hoja, reproduce los 5 acuerdos del uso responsable de IA. Al lado de cada uno, palomita ✓ si te comprometes. Si dudas con alguno, conversa con el profe.\\
Abstracción & \textbf{Paso 3 --- 5 preguntas-radar.} Lista numerada de 5 preguntas tuyas sobre IA. Sugerencias para inspirarte: \emph{``¿Cómo funciona realmente ChatGPT por dentro?''}, \emph{``¿Es seguro usar IA para tareas del colegio?''}, \emph{``¿La IA va a quitarme el trabajo cuando sea adulto?''}, \emph{``¿Cómo se cuáles son las mejores IAs?''}, \emph{``¿Qué es prompting?''}.\\
Algoritmo de armado & \textbf{Paso 4 --- Mini-investigación del consejero.} En la siguiente hoja, subtítulo: \textbf{``Mini-investigación: el consejero del barrio''}. Esta tarea la haces en casa: pregunta a una persona mayor de tu familia (mamá, papá, abuelos, tíos): \emph{``¿Quién era el consejero del barrio en su tiempo? ¿Qué tipo de consejos daba?''}. Apunta mínimo \textbf{5 datos concretos}: nombre del consejero, barrio o vereda, tipo de consejos, qué cualidades tenía, una historia concreta.\\
\end{tabularx}
\normalsize

\begin{drawbox}{milcMagenta}{Antes de cerrar la S1 del P3}
\checkbox El índice del P3 está con 10 sesiones. \\[1mm]
\checkbox Los 5 acuerdos están firmados. \\[1mm]
\checkbox Tengo 5 preguntas-radar reales. \\[1mm]
\checkbox Tengo plan para investigar al consejero tradicional. \\[1mm]
\checkbox La sesión 1 tiene encabezado y registro. \\[1mm]
\checkbox Está firmado con nombre y fecha.
\end{drawbox}

\begin{softbox}{milcMagenta}{milcGris}{Plantilla de trabajo}
\textbf{Estructura.} \textit{Hoja 1:} título P3 + índice de 10 sesiones. \textit{Hoja 2:} 5 acuerdos firmados + 5 preguntas. \textit{Hoja 3:} mini-investigación del consejero (5+ datos). \textit{Cuaderno principal:} S1 con encabezado fijo.
\end{softbox}

\begin{infoband}{milcMagenta}


\textbf{Lo que va al cuaderno (Actividad 3 · CREA).} Título: «Actividad 3 --- Cuaderno inaugural G7-P3». \textbf{Formato:} 3 hojas iniciales + S1 registrada. \textbf{Extensión:} 3 hojas + sesión. \textbf{Sabes que terminaste cuando} tu cuaderno está listo para empezar la S2.
\end{infoband}

\puente{El producto es el cuaderno P3 + 5 acuerdos firmados + investigación.}

\begin{titlebox}{milcMostaza}
Producto de la sesión
\end{titlebox}
\textbf{Cuaderno P3 inaugural + 5 acuerdos + 5 preguntas + investigación del consejero}.
\begin{softbox}{milcMostaza}{milcGris}{Criterios de calidad}
(1) Índice del P3 con 10 sesiones. (2) 5 acuerdos firmados con palomita. (3) 5 preguntas-radar personales sobre IA. (4) Plan claro de mini-investigación. (5) S1 con encabezado fijo y registro.
\end{softbox}

\puente{Antes de salir, verifica que tu cuaderno tiene los 5 elementos del periodo.}

\begin{titlebox}{milcVino}
Fase 4 · Evaluación liberadora — cinco dimensiones
\end{titlebox}

\begin{softbox}{milcVino}{milcGris}{Sentido de la evaluación}
Evaluar no es solo revisar una nota. Es reconocer qué aprendí, qué cambió en mí, qué emociones manejé, cómo me relaciono con la ciudad y la comunidad, y qué le dejaré a quien viene detrás.
\end{softbox}

\textbf{Mi reflexión liberadora en cinco dimensiones.} Completa cada frase iniciada:

\small
\begin{tabularx}{\textwidth}{>{\bfseries\RaggedRight\arraybackslash}p{3.4cm}Y>{\RaggedRight\arraybackslash}p{4.6cm}}
\rowcolor{milcVino}\color{white}Dimensión & \color{white}\bfseries Mi respuesta (frame starter) & \color{white}\bfseries Algo que descubrí\\
1. Personal & Lo que más cambió en mí fue \linead{6.5cm} & \linead{4.2cm}\\
2. Emocional & La emoción más fuerte fue \linead{2.5cm} y la regulé \linead{3cm} & \linead{4.2cm}\\
3. Ciudadana & En democracia esto importa porque \linead{6cm} & \linead{4.2cm}\\
4. Local (Cartago) & En mi barrio sería útil para \linead{6.5cm} & \linead{4.2cm}\\
5. Intergeneracional & A mi abuela/abuelo le diría \linead{6.5cm} & \linead{4.2cm}\\
\end{tabularx}
\normalsize

\begin{infoband}{milcVino}
\textbf{Para la próxima sustentación.} Vuelve a esta tabla en seis meses. Verás que las respuestas cambian — no porque lo aprendido cambie, sino porque tú cambias. La evaluación liberadora es una conversación contigo a través del tiempo.
\end{infoband}


\puente{Cerramos con 3 voces sobre por qué pensar bien sobre la IA hoy decide quién serás en 10 años.}

\begin{titlebox}{milcGradeDark}
Triángulo de pensamiento · cierre filosófico
\end{titlebox}

\begin{softbox}{milcVino}{milcGris}{Dussel · lente del nosotros}
\textit{````El consejero del barrio sigue siendo necesario, aunque cambie de forma. Hoy es humano + IA juntos.''''}

\textbf{El consejero del barrio no desaparece: cambia de forma}. Tu generación va a usar IA como herramienta de consejo, sí, pero \emph{el consejero humano sigue valiendo}: el profe, el papá, la abuela. \emph{La IA no reemplaza la sabiduría humana --- la complementa}. Doña Mercedes seguiría siendo útil aunque existiera ChatGPT. \emph{Lo importante es no perder el consejero humano por usar el digital}.

\textbf{¿A quién considero mi ``consejero del barrio'' moderno? ¿Cómo combino consejo humano + IA con criterio?}
\end{softbox}
\linead{\linewidth}\\[1mm]
\linead{\linewidth}

\begin{softbox}{milcMostaza}{milcGris}{Estoicismo · Marco Aurelio · cuidado interior}
\textit{````El que pide consejo no es débil; el que confía sin verificar sí lo es. La IA es consejo nuevo: trátalo con la sabiduría antigua.''''}

Marco Aurelio \emph{pedía consejo todos los días} a su Senado, a sus asesores, a sus generales. \textbf{Pero después él decidía}. No delegaba la decisión final. La IA es un consejero nuevo en tu lista. Trátalo con la misma sabiduría: \emph{escucha + verifica + decide tú}. Esa disciplina antigua es lo que distingue a quien usa IA con criterio de quien se convierte en su esclavo.

\textbf{Cuando uso IA, ¿estoy decidiendo yo después de escuchar, o estoy delegando mi decisión a ella?}
\end{softbox}
\linead{\linewidth}\\[1mm]
\linead{\linewidth}

\begin{softbox}{milcTurquesa}{milcGris}{Floridi · lente de la infoesfera}
\textit{````La generación que aprende IA en colegio decide cómo será la humanidad en 30 años. No es exageración: es realidad.''''}

\textbf{Tu generación es la primera que crece con IA conversacional accesible}. ChatGPT salió a fines de 2022. Tú entras a grado 7 en 2026 con la IA ya como parte del paisaje. \emph{Las decisiones que tomas hoy sobre cómo usar IA forman los hábitos que cargarás 50 años}. Por eso este periodo importa más que cualquier otro: \emph{no es solo ``aprender a usar la app''; es decidir qué tipo de adulto digital vas a ser}.

\textbf{Si mi yo de 30 años pudiera verme hoy usando IA, ¿estaría orgulloso de los hábitos que estoy formando?}
\end{softbox}
\linead{\linewidth}\\[1mm]
\linead{\linewidth}

\begin{titlebox}{milcVino}
Compromiso final
\end{titlebox}

\small
\begin{tabularx}{\textwidth}{>{\RaggedRight\arraybackslash}p{3.0cm}YYY}
\rowcolor{milcVino}\color{white}\bfseries Criterio & \color{white}\bfseries Lo logré & \color{white}\bfseries En proceso & \color{white}\bfseries Necesito apoyo\\
Comprensión del tema & Explico ideas centrales con mis palabras. & Reconozco ideas pero falta precisión. & Necesito repasar conceptos base.\\
Pensamiento computacional & Apliqué los 4 pilares al diseñar mi producto. & Apliqué algunos pilares. & Necesito releer la fase 2.\\
Aplicación y producto & Mi producto cumple los criterios y se puede revisar. & Producto funcional con ajustes pendientes. & Aún en construcción.\\
Reflexión 5 dimensiones & Respondí cada dimensión con honestidad. & Respondí algunas con detalle. & Mis respuestas son muy generales.\\
Triángulo de pensamiento & Conecté las 3 voces con mi trabajo real. & Conecté al menos una. & No relaciono las voces con mi proceso.\\
\end{tabularx}
\normalsize

\textbf{Hoy entendí que:}\\[1mm]
\linead{\linewidth}\\[1mm]
\linead{\linewidth}

\textbf{Mi compromiso para la próxima sesión es:}\\[1mm]
\linead{\linewidth}\\[1mm]
\linead{\linewidth}

\begin{softbox}{milcMostaza}{milcGris}{Cita de despedida · Marco Aurelio}
\textit{``El obstáculo en el camino se vuelve el camino. Lo que estorba al camino se vuelve camino.''}\\[1mm]
Cada error de hoy, bien observado, se convierte en pieza del camino que pisarás mañana.
\end{softbox}

\small
\begin{tabularx}{\textwidth}{>{\bfseries}p{3.6cm}Y}
\rowcolor{milcGradeDark!12}Nombre completo: & \linead{10cm}\\
Fecha: & \linead{4cm} \quad Curso/grupo: \linead{4cm}\\
Mi firma: & \linead{9cm}\\
\end{tabularx}
\normalsize

\begin{infoband}{milcGradeDark}
\textbf{Para la próxima sesión.} Esta semana, termina la mini-investigación del consejero (pregunta en casa). La próxima clase es S2: \emph{¿Qué es la IA?}. Vas a aprender la definición técnica, la breve historia y 10 casos cotidianos donde ya está cambiando tu vida sin que lo notes.
\end{infoband}

\end{document}
